\documentclass{article}
\usepackage{PRIMEarxiv} % Core package for the PrimeAI Template
\usepackage{amsmath, amssymb, amsthm, bm, dcolumn} % Add amsthm here for the proof environment
\usepackage[numbers,sort&compress]{natbib} % Natbib for citations
\usepackage{graphicx} % For high-quality images
\usepackage{multicol} % Multi-column figures/tables
\usepackage{hyperref} % Hyperlinks
\usepackage{listings} % Code listings
\usepackage{authblk} % For structured affiliations
% Define theorem style (optional)
\theoremstyle{plain}
\newtheorem{theorem}{Theorem}
\renewcommand{\qedsymbol}{}

% Custom commands for primes and qubit encoding
\newcommand{\primeQubit}[1]{|p_{#1}\rangle}
\newcommand{\primeGate}[1]{U_{p_{#1}}}

\usepackage{wrapfig}
\usepackage[pscoord]{eso-pic}
\usepackage[fulladjust]{marginnote}
\reversemarginpar

% Typesetting improvements without footnote patching
\usepackage[protrusion=true, expansion=true, tracking=false]{microtype}
\microtypecontext{spacing=nonfrench}

% Line numbers
\usepackage[right]{lineno}

% Text layout - adjust as needed
\raggedright
\setlength{\parindent}{0.5cm}
\textwidth 5.25in 
\textheight 8.75in

% Set double spacing
\usepackage{setspace} 
\doublespacing

% Adjust width for specific content
\usepackage{changepage}

% Adjust caption style
\usepackage[aboveskip=1pt,labelfont=bf,labelsep=period,singlelinecheck=off]{caption}

% Remove brackets from references
\makeatletter
\renewcommand{\@biblabel}[1]{\quad#1.}
\makeatother

% Header, footer, and page numbers
\usepackage{lastpage,fancyhdr}
\pagestyle{fancy}
\fancyhf{}
\fancyhead[L]{Preprint - PrimeAI Enhanced Template}
\fancyfoot[C]{\scriptsize Multiplicity Theory © 2024 Ryan Van Gelder - Citizen Gardens \\ Licensed Under MIT and CC BY-NC-SA 4.0.}
\fancyfoot[R]{Page \thepage\ of \pageref{LastPage}}
\renewcommand{\footrule}{\hrule height 2pt \vspace{2mm}}

\begin{document}

\title{Integrating Prime Encoding into Minkowski Distance}
\author{Ryan O. Van Gelder}
\affil{Citizen Gardens - The Foundation of Multiplicity \\ \texttt{info@citizengardens.org}}
\date{}
\maketitle

\section*{1. Minkowski Distance Formula}
The Minkowski distance between two points \( \mathbf{x} = (x_1, x_2, \dots, x_n) \) and \( \mathbf{y} = (y_1, y_2, \dots, y_n) \) in \( \mathbb{R}^n \) is defined as:
\begin{equation}
d(\mathbf{x}, \mathbf{y}) = \left( \sum_{i=1}^n |x_i - y_i|^p \right)^{1/p},
\end{equation}
where \( p \geq 1 \) determines the order of the distance metric.

\section*{2. Prime Encoding Scheme}
To encode coordinates \( x_i \) and \( y_i \) with prime numbers:
\begin{itemize}
    \item Assign a unique prime number \( p_i \) to each dimension \( i \).
    \item Encode the coordinates as:
    \begin{equation}
    \phi(x_i) = p_i \cdot x_i, \quad \phi(y_i) = p_i \cdot y_i.
    \end{equation}
\end{itemize}

\subsection*{Encoded Difference}
The difference in encoded coordinates is:
\begin{equation}
\Delta_i = \phi(x_i) - \phi(y_i) = p_i \cdot (x_i - y_i).
\end{equation}

\section*{3. Prime-Encoded Minkowski Distance}
Substituting the encoded difference into the Minkowski formula:
\begin{equation}
d_{\text{prime}}(\mathbf{x}, \mathbf{y}) = \left( \sum_{i=1}^n |p_i \cdot (x_i - y_i)|^p \right)^{1/p}.
\end{equation}

\section*{4. Properties of Prime Encoding}
\subsection*{a. Uniqueness}
Primes ensure each dimension is independent:
\begin{equation}
\text{If } \phi(x_i) \neq \phi(y_i), \quad \text{then } x_i \neq y_i \text{ for all } i.
\end{equation}

\subsection*{b. Error Detection}
Errors in computations can be detected as:
\begin{equation}
\text{Error} = \phi(x_i) \mod p_i \neq x_i \cdot p_i.
\end{equation}

\subsection*{c. Compactness}
The total prime product \( P = \prod_{i=1}^n p_i \) provides a compact representation of the coordinate system.

\section*{5. Recursive Feedback for Adaptive Encoding}
Recursive feedback loops adjust prime mappings dynamically:
\begin{equation}
p_i(t+1) = p_i(t) \cdot F(t),
\end{equation}
where \( F(t) \) is a feedback function dependent on system dynamics.

\subsection*{Dynamic Update}
The encoded Minkowski distance evolves as:
\begin{equation}
d_{\text{prime}}^{(t+1)}(\mathbf{x}, \mathbf{y}) = \left( \sum_{i=1}^n |p_i(t+1) \cdot (x_i - y_i)|^p \right)^{1/p}.
\end{equation}

\section*{6. Integration with Multiplicity Framework}
The time-dependent multiplicity operator \( M(t, \psi(t)) \) incorporates prime-encoded Minkowski metrics:
\begin{equation}
M(t, \psi(t)) = \sum_{i=1}^n p_i(t) \cdot \psi_i(t),
\end{equation}
where \( \psi_i(t) \) is the state vector.

\subsection*{Coupled with Tensor Networks}
The coupling tensor \( T(t) \) models higher-dimensional interactions:
\begin{equation}
T_{ij}(t) = p_i(t) \cdot p_j(t).
\end{equation}

\subsection*{Prime-Encoded Multiplicity Equation}
\begin{equation}
H(t) \ni \psi(t) \to \left( M(t, \psi(t)) \cdot T(t, \psi(t)) + f(t, \psi(t)) \right) = \lambda(t) \psi(t).
\end{equation}

\section*{7. Stochastic Extensions}
Introduce noise \( \epsilon_i(t) \) for real-world adaptability:
\begin{equation}
d_{\text{prime}}(\mathbf{x}, \mathbf{y}) = \left( \sum_{i=1}^n |(p_i \cdot (x_i - y_i) + \epsilon_i(t))|^p \right)^{1/p}.
\end{equation}

\section*{8. Applications}
\begin{itemize}
    \item \textbf{Quantum Systems:} Prime-encoded Minkowski distances support quantum superposition and coherence:
    \begin{equation}
    \psi(t) = \sum_{i=1}^n p_i(t) \cdot \alpha_i \cdot e^{i \theta_i}.
    \end{equation}
    \item \textbf{Error Correction:} Errors are corrected using:
    \begin{equation}
    \phi_{\text{corrected}}(x_i) = \frac{\phi(x_i)}{\gcd(\phi(x_i), E)},
    \end{equation}
    where \( E \) is the error factor.
\end{itemize}

\section*{Conclusion}
Prime encoding of Minkowski distance enhances precision, scalability, and adaptability while aligning with Multiplicity Theory. This framework supports quantum systems, AI, and high-dimensional computations.


\end{document}
